% Standalone problem document, generated from the adjacent README.md.
% Compile with XeLaTeX. Mathematical content is maintained in Markdown.
\documentclass[11pt,a4paper]{article}
\usepackage[fontsize=10.5pt]{scrextend}
\usepackage[margin=22mm,headheight=15pt,headsep=9mm,footskip=11mm]{geometry}
\usepackage{fontspec}
\setmainfont{texgyrepagella}[Extension=.otf,UprightFont=*-regular,BoldFont=*-bold,ItalicFont=*-italic,BoldItalicFont=*-bolditalic]
\setsansfont{texgyreheros}[Extension=.otf,UprightFont=*-regular,BoldFont=*-bold,ItalicFont=*-italic,BoldItalicFont=*-bolditalic]
\setmonofont{texgyrecursor}[Extension=.otf,UprightFont=*-regular,BoldFont=*-bold,ItalicFont=*-italic,BoldItalicFont=*-bolditalic,Scale=MatchLowercase]
\usepackage{amsmath,amssymb}
\usepackage{unicode-math}
\setmathfont{texgyrepagella-math.otf}
\usepackage{xcolor,microtype,enumitem,needspace,fancyhdr}
\usepackage{bookmark}
\definecolor{ink}{HTML}{17344B}
\definecolor{accent}{HTML}{197D80}
\definecolor{muted}{HTML}{596773}
\hypersetup{colorlinks=true,linkcolor=accent,urlcolor=accent,pdftitle={SP-02 - The
sharp off-diagonal threshold for spectral-subspace
rotation},pdfauthor={Open Problems in Numerical Linear Algebra}}
\urlstyle{same}
\setlength{\parindent}{0pt}
\setlength{\parskip}{5pt plus 1pt minus 1pt}
\linespread{1.04}
\setlength{\emergencystretch}{3em}
\widowpenalty=10000
\clubpenalty=10000
\displaywidowpenalty=10000
\allowdisplaybreaks[1]
\setlist{nosep,leftmargin=1.5em}
\providecommand{\tightlist}{\setlength{\itemsep}{0pt}\setlength{\parskip}{0pt}}
\providecommand{\pandocbounded}[1]{#1}
\pagestyle{fancy}
\fancyhf{}
\fancyhead[L]{\sffamily\footnotesize\color{muted}OPEN PROBLEMS IN NUMERICAL LINEAR ALGEBRA}
\fancyhead[R]{\sffamily\bfseries\footnotesize\color{accent}SP-02}
\fancyfoot[L]{\parbox[b]{0.9\textwidth}{\sffamily\fontsize{8}{10}\selectfont\color{muted}Editorial ratings; a bounded search cannot certify that a problem remains unsolved.\\Literature check: 2026-09-10}}
\fancyfoot[R]{\sffamily\footnotesize\color{muted}\thepage}
\renewcommand{\headrulewidth}{0.3pt}
\renewcommand{\headrule}{\hbox to\headwidth{\color{accent}\leaders\hrule height \headrulewidth\hfill}}
\makeatletter
\renewcommand{\section}{\Needspace{5\baselineskip}\@startsection{section}{1}{0pt}{12pt plus 2pt minus 2pt}{4pt}{\sffamily\bfseries\large\color{ink}}}
\renewcommand{\subsection}{\Needspace{4\baselineskip}\@startsection{subsection}{2}{0pt}{9pt plus 1pt minus 1pt}{3pt}{\sffamily\bfseries\normalsize\color{ink}}}
\makeatother
\setcounter{secnumdepth}{0}
\begin{document}
{\sffamily\small\color{muted}Eigenvalues and inverse problems\par}
\vspace{3pt}
{\raggedright\hyphenpenalty=10000\exhyphenpenalty=10000\sffamily\bfseries\fontsize{21}{24}\selectfont\color{ink}The
sharp off-diagonal threshold for spectral-subspace rotation\par}
\vspace{7pt}
{\sffamily\small\textbf{Difficulty:} challenging\quad\textbf{Importance:} interesting
to the community\par}
{\sffamily\small\bfseries\color{accent}Status: Open\par}
\vspace{4pt}
{\color{accent}\hrule height 0.8pt}
\vspace{6pt}
\textbf{Rating rationale:} Challenging reflects an optimal off-diagonal
perturbation bound for arbitrary separated spectral sets; community
impact is a sharp guarantee for invariant-subspace stability.

Let \(A\) be a self-adjoint, possibly unbounded operator on a separable
complex Hilbert space, with \(\sigma(A)=\sigma\cup\Sigma\) for nonempty
closed sets satisfying \(d=\operatorname{dist}(\sigma,\Sigma)>0\). Let
\(P=E_A(\sigma)\). Suppose a bounded self-adjoint perturbation \(V\) is
off-diagonal relative to this decomposition: \[
PVP=0,\qquad (I-P)V(I-P)=0.
\] Writing \(O_r(S)=\{x\in\mathbb R:\operatorname{dist}(x,S)<r\}\) and
\(Q=E_{A+V}(O_{d/2}(\sigma))\), does \[
\|V\|<\frac{\sqrt3}{2}d\quad\Longrightarrow\quad\|P-Q\|<1
\] hold universally? Here \(E_B\) is the spectral projection measure of
\(B\), and norms are operator norms. No ordering of the two spectral
sets is assumed.

The source retains the full operator setting, encompassing Hermitian
block matrices in invariant-subspace computations. The threshold
\(\sqrt3/2\) already guarantees that the perturbed spectral components
stay separated; the unresolved assertion controls their maximal subspace
angle. Off-diagonal structure makes this a different threshold problem
from the general-perturbation constant \(1/2\).

\subsection{References}\label{references}

A. Seelmann, \emph{Notes on the subspace perturbation problem for
off-diagonal perturbations}, Proceedings AMS 144 (2016), 3825--3832, §1,
(1.1)--(1.4), and Theorem 2.5
(\href{https://arxiv.org/pdf/1412.6294}{primary manuscript};
\href{https://doi.org/10.1090/proc/13118}{journal}). A. Seelmann,
\emph{Unifying the treatment of indefinite and semidefinite
perturbations in the subspace perturbation problem}, Operators and
Matrices 15 (2021), 1181--1188
(\href{https://files.ele-math.com/articles/oam-15-74.pdf}{primary
paper}), treats a different extension of the perturbation hypotheses.

\subsection{Status check --- 2026-09-08}\label{status-check-2026-09-08}

The 2016 paper explicitly conjectures \(c_{\mathrm{opt-off}}=\sqrt3/2\)
and proves \(c_{\mathrm{opt-off}}>0.6940725\). The 2021 extension does
not attain this off-diagonal endpoint. Searches for ``off-diagonal
subspace perturbation'', ``0.6940725'', ``sqrt 3'', ``optimal
constant'', ``conjecture'', and 2024--2026 found no resolution. This is
a historical-source entry with a bounded later-literature check, not a
certified claim of present openness.

\subsection{Audit update --- 2026-09-10}\label{audit-update-2026-09-10}

Rechecked \href{https://arxiv.org/pdf/1412.6294}{Seelmann's account} and
the \href{https://files.ele-math.com/articles/oam-15-74.pdf}{2021
follow-up}. The proved sufficient constant above 0.694 remains below the
conjectured threshold \(\sqrt3/2\). Searches for subsequent off-diagonal
optimal-threshold results found no general proof or counterexample. A
smaller sufficient constant is progress but not a proved subfamily of
the optimal-constant determination.
\end{document}
